\section*{Chapter \ref{ch:m3:maximal-extractable-value} --- Maximal Extractable Value}

\begin{solution}{exo:m3:maximal-extractable-value:1}
Any front-run lowers the victim's output below the quote, so the swap would revert and the attacker would be left with
its front-run position. Wallets allow some tolerance because the pool can move for innocent reasons between submission and
inclusion (other trades, arbitrage), and a zero-tolerance swap would often fail.
\end{solution}

\begin{solution}{exo:m3:maximal-extractable-value:2}
The liquidation after an \omterm{def:m3:blockchains-for-traders:contract}{oracle} update is a back-run; the arbitrage after a large swap is a back-run; the purchase before a
large purchase is a front-run.
\end{solution}

\begin{solution}{exo:m3:maximal-extractable-value:3}
To the builder: the proposer cannot see or change the block's transactions before committing to propose it. To the
proposer: the bid it accepts is valid and will be paid if it proposes the block.
\end{solution}

\begin{solution}{exo:m3:maximal-extractable-value:4}
Both double: USD~5\,067 to 10\,130 of gross gain, and 1.56 to 3.12 ether of loss.
\end{solution}

\begin{solution}{exo:m3:maximal-extractable-value:5}
$10\,129.57 \times 10\% - 20 = $ USD~992.96.
\end{solution}

\begin{solution}{exo:m3:maximal-extractable-value:6}
Channels send flow to the builders they trust or are paid by; the builders with most private flow can build the most valuable
blocks, win more auctions, and attract more flow, a feedback loop.
\end{solution}

\begin{solution}{exo:m3:maximal-extractable-value:7}
A rise of about 0.54\% (to 3\,016.20): the arbitrage buys about USD~17\,900 of ether for a profit of about USD~21. Below the
pool's 0.30\% fee there is no arbitrage at all.
\end{solution}

\begin{solution}{exo:m3:maximal-extractable-value:8}
The channel prevents sandwiches but the swap still moves the pool, whose back-run value is shared or kept by the channel and
builders; \omterm{def:m3:blockchains-for-traders:gas}{priority fees} and the channel's own terms are costs; and the trader trusts the channel and its builders not to use
the flow. MEV exposure is reduced and moved, not removed.
\end{solution}

\begin{solution}{pb:m3:maximal-extractable-value:1}
\textbf{1.} 311.62 ether, an average of 3\,209.03 dollars: about 7\% above the pool's price, of which 0.30\% is the fee and
the rest the swap's own impact. \textbf{2.} 310.06 and 308.50 ether. \textbf{3.} USD~38\,912 and 78\,119. \textbf{4.} Its
gain rises with the front-run up to the point where the victim would revert. \textbf{5.} Nothing: no front-run is feasible.
\textbf{6.} USD~5\,067 and 10\,130. \textbf{7.} 1.56 and 3.12 ether, about USD~4\,674 and 9\,349 at 3\,000. \textbf{8.}
USD~5\,046.72 and 10\,109.57 with no payment; USD~486.67 and 992.96 paying 90\%. \textbf{9.} The builder, and through its
bid the proposing \omterm{def:m3:blockchains-for-traders:pow}{validator}. \textbf{10.} Every \omterm{def:m3:maximal-extractable-value:mev}{searcher} who sees the same opportunity bids for it; the highest bid approaches
the whole value. \textbf{11.} Just above the swap's expected slippage from other trades before inclusion, a few basis points, or
send it privately; 0.5\% is a gift of up to USD~4\,700. \textbf{12.} Each piece moves the pool less, so each tolerance can be
tighter and each sandwich smaller; the trader bears price risk across blocks and ten times the gas. \textbf{13.} It hides the
swap from public \omterm{def:m3:maximal-extractable-value:mev}{searchers} and can refund back-run value; it does not remove the price impact or the trust in the channel.
\textbf{14.} Solvers or fillers compete to fill at one price, and within a batch the order does not matter; the trader's worst
price is still its signed limit. \textbf{15.} Look in the block for the same address buying the same \omterm{def:m3:blockchains-for-traders:key}{token} just before and
selling just after, in the same pool (\texttt{find\_sandwiches}), and compare the fill with the quote. \textbf{16.} It exploits
knowledge of a pending order to trade ahead of it, which is the substance of \omterm{def:m3:maximal-extractable-value:front}{front-running}, but the attacker owes the victim no
duty as a broker does; whether it is illegal depends on jurisdiction and facts, and it is contested. \textbf{17.} They trade on
information that is public after the transaction and do not worsen the user's price; they correct prices. \textbf{18.} That the
separation of roles depends on software that can have flaws, and that participants who control \omterm{def:m3:blockchains-for-traders:pow}{validators} can exploit them;
trust has moved to relays and their code. \textbf{19.} \emph{Named result:} USD~5\,046.72 net and 1.56 ether lost at 0.5\%;
USD~10\,109.57 net and 3.12 ether lost at 1\%. \textbf{20.} A written option on the swap's price, exercisable by whoever orders
the block.
\end{solution}

\begin{solution}{iq:m3:maximal-extractable-value:1}
A buy before and a sell after a victim's swap in the same pool, bounded by the victim's \omterm{def:m3:maximal-extractable-value:sandwich}{slippage tolerance}. Protect by tight
tolerances, private channels, splitting large swaps, and intent or batch mechanisms.
\iqlookfor{tolerance as the bound, and private flow.}
\end{solution}

\begin{solution}{iq:m3:maximal-extractable-value:2}
Builders assemble blocks and bid for inclusion; proposers pick the highest bid without seeing contents; relays hold the block
contents until the proposer commits (a commit-and-reveal) and check bids, so that neither side can cheat the other.
\iqlookfor{the commit and reveal and who trusts whom.}
\end{solution}

\begin{solution}{iq:m3:maximal-extractable-value:3}
The victim's output falls with the front-run; set the output after a front-run $a$ equal to $(1 - \tau)$ times the quote and
solve for $a$ (in closed form for the constant product, or by binary search in integer arithmetic).
\iqlookfor{monotonicity and the binding constraint.}
\end{solution}

\begin{solution}{iq:m3:maximal-extractable-value:4}
It is the same flow seen from two sides: the arbitrageur's profit from trading the stale pool to the external price is the pool's
loss to rebalancing, net of the fee.
\iqlookfor{the arbitrageur's gain equals the pool's loss.}
\end{solution}

\begin{solution}{iq:m3:maximal-extractable-value:5}
Operational: bugs in bundles that lose money, key management, dependence on relays and builders, reverted transactions and gas.
Legal: strategies that harm users (sandwiches) may be treated as fraud or market manipulation; \omterm{def:m3:blockchains-for-traders:pow}{validators} or exploits of
infrastructure bring criminal exposure, as the 2023 relay case shows.
\iqlookfor{separation of benign \omterm{def:m3:maximal-extractable-value:front}{back-running} from harmful strategies.}
\end{solution}

\begin{solution}{iq:m3:maximal-extractable-value:6}
Users send transactions to a trusted node that reveals only hints (pool, direction, not size or limit); \omterm{def:m3:maximal-extractable-value:mev}{searchers} bid for the
right to back-run with bundles the node simulates; the node forwards the best bundle with the condition that the user is paid a
stated share, and only back-runs are accepted. Harden with sealed bids, several builders and audits.
\iqlookfor{hints not orders, back-runs only, and payment conditions.}
\end{solution}
